% ============================================================
% Chapter 9 — Results
% ============================================================
\chapter{Results}
\label{ch:results}

\roadmap{This chapter reports the user-study results. It is written as
complete scaffolding: every empirical value is a \todo{} placeholder to be
filled once the study data are frozen, but the section structure, table
skeletons, and figure placeholders are final, and each is tied to the
method that produces it in \cref{ch:methodology}. Where a result depends on
an externally defined grouping (notably the phase contrast of item 16),
the operationalization is restated at the point of use.}

\section{Participants and data yield}
\label{sec:results-participants}

The study enrolled \todo{N} participants
(\todo{demographic summary: age range, prior experience, recruitment}).
Data collection yielded \todo{S} learning sessions totaling \todo{H} hours
of instrumented interaction. Of these, \todo{S_sr} sessions post-date the
self-report probe (2026-07-01) and are usable for the item-15 validation;
the remaining sessions predate it (\cref{sec:methodology-selfreport}).
\Cref{tab:results-yield} summarizes per-channel data yield, including the
fraction of frames with a detected face, which bounds the facial analyses
\citep{bosch2015}.

\begin{table}[htbp]
\centering
\small
\caption{Per-channel data yield.}
\label{tab:results-yield}
\begin{tabular}{@{}p{0.36\linewidth}p{0.24\linewidth}p{0.32\linewidth}@{}}
\toprule
Channel & Volume & Coverage / quality \\
\midrule
Learning sessions & \todo{S} & \todo{mean duration} \\
Webcam emotion frames & \todo{rows} & \todo{\% frames with face} \\
Action-unit frames & \todo{rows} & \todo{\% frames with face} \\
Gaze samples & \todo{rows} & \todo{mean confidence; \% above 0.5} \\
Self-report responses & \todo{rows} & \todo{responses per session} \\
Chat / dialogue utterances & \todo{rows} & \todo{EF detections} \\
\bottomrule
\end{tabular}
\end{table}

\section{Platform quality attainment}
\label{sec:results-quality}

Against the eighteen-item accounting of \cref{sec:scope-table}, the
delivered platform attains \todo{count} items as delivered or exceeded,
\todo{count} as partial or reframed, and \todo{count} as unverifiable from
code, exactly as \cref{tab:scope-18} records; no status is upgraded here.
\todo{If a rubric-based quality score was collected (e.g., a usability or
system-quality questionnaire), report it here with its instrument and
scale.}

\section{Cross-modal validation (RQ1)}
\label{sec:results-validation}

Following \cref{sec:methodology-validation}, engagement estimates were
compared with gaze-on-AOI on-task labels over \todo{window length} windows.
The continuous association was $r = \todo{value}$ (\todo{95\% CI};
$p = \todo{value}$; $n = \todo{windows}$ windows across \todo{participants}
participants). Binarized agreement was $\kappa = \todo{value}$
(\todo{interpretation band}). \Cref{fig:validation-confusion} shows the
self-report-by-detected-state confusion matrix, and
\cref{tab:results-validation} reports sensitivity and specificity per
state.

\begin{figure}[htbp]
\centering
\figplaceholder{Validation confusion matrix}{Rows: five self-reported
states; columns: four detected states (plus unresolved). Cells: normalized
frequency}
\caption[Self-report vs.\ detected-state confusion matrix]{Agreement
between five-option self-report and system-detected states. See
\cref{sec:methodology-validation}.}
\label{fig:validation-confusion}
\end{figure}

\begin{table}[htbp]
\centering
\small
\caption{Detection agreement per state.}
\label{tab:results-validation}
\begin{tabular}{@{}lccc@{}}
\toprule
State & Sensitivity & Specificity & Base rate \\
\midrule
Engagement & \todo{} & \todo{} & \todo{} \\
Boredom & \todo{} & \todo{} & \todo{} \\
Confusion & \todo{} & \todo{} & \todo{} \\
Frustration & \todo{} & \todo{} & \todo{} \\
\bottomrule
\end{tabular}
\end{table}

\todo{State whether the facial engagement estimate or the AU-feature
variant was used as primary, and report the other as a robustness check.}

\section{Executive-function detection distributions}
\label{sec:results-ef}

After excluding error rows, \todo{count} EF detections were analyzed.
\Cref{fig:ef-distribution} shows detection rates per construct; mean
confidence per construct is reported in \cref{tab:results-ef}, with the two
low-feasibility constructs (cognitive flexibility, working memory) flagged
as the implementation flags them.

\begin{figure}[htbp]
\centering
\figplaceholder{EF detection distribution}{Bar chart of detection
count/rate for each of the nine constructs, error rows excluded}
\caption[EF detection distribution]{Detection distribution across the nine
constructs.}
\label{fig:ef-distribution}
\end{figure}

\begin{table}[htbp]
\centering
\small
\caption{EF detections by construct (error rows excluded).}
\label{tab:results-ef}
\begin{tabular}{@{}lcc@{}}
\toprule
Construct & Detections & Mean confidence \\
\midrule
\todo{one row per construct} & \todo{} & \todo{} \\
\bottomrule
\end{tabular}
\end{table}

\section{Learning gain (RQ2)}
\label{sec:results-gain}

Using the pre/post design of \cref{sec:methodology-gain} with a
\todo{first/best/last} attempt policy, \todo{n} learners completed both
assessments. Mean pre score was \todo{value}, mean post score
\todo{value}; the paired comparison gave \todo{$t$ or $W$} $= \todo{}$,
$p = \todo{}$, with effect size \todo{$d$ or rank-biserial} $= \todo{}$.
Mean normalized gain was $\langle g \rangle = \todo{}$.
\Cref{fig:gain-slopegraph} shows per-learner pre$\rightarrow$post
trajectories, and \cref{tab:results-gain} reports per-knowledge-component
gains.

\begin{figure}[htbp]
\centering
\figplaceholder{Learning-gain slopegraph}{Per-learner pre$\rightarrow$post
lines with the normalized-gain distribution alongside}
\caption[Learning gain]{Per-learner pre/post trajectories and normalized
gain.}
\label{fig:gain-slopegraph}
\end{figure}

\begin{table}[htbp]
\centering
\small
\caption{Learning gain by knowledge component.}
\label{tab:results-gain}
\begin{tabular}{@{}lccc@{}}
\toprule
Knowledge component & Pre & Post & $\langle g \rangle$ \\
\midrule
\todo{one row per KC} & \todo{} & \todo{} & \todo{} \\
\bottomrule
\end{tabular}
\end{table}

\section{Intervention engagement}
\label{sec:results-interventions}

\Cref{tab:results-interventions} reports, per strategy, how often
interventions were initiated, completed, and dismissed
(\cref{sec:methodology-analyses}, analysis 1), and the association between
chatbot strategy suggestions and subsequent learner-initiated use
(\cref{sec:interventions-suggest}).

\begin{table}[htbp]
\centering
\small
\caption{Intervention initiation and completion by strategy.}
\label{tab:results-interventions}
\begin{tabular}{@{}lccc@{}}
\toprule
Strategy & Initiated & Completed & Dismissed \\
\midrule
Practice testing & \todo{} & \todo{} & \todo{} \\
Interrogative elaboration & \todo{} & \todo{} & \todo{} \\
Stepwise learning & \todo{} & \todo{} & \todo{} \\
Distributed practice & \todo{} & \todo{} & \todo{} \\
\bottomrule
\end{tabular}
\end{table}

\section{Group comparison for proposal items 16--17}
\label{sec:results-groupcompare}

\emph{Operationalization.} Because the database encodes no phase marker
(\cref{sec:methodology-stats}), the two ``phases'' are operationalized as
the pre-study and post-study conditions. Under this definition, the
dependent variable \todo{name the DV} was compared between conditions.
A Shapiro--Wilk test \citep{shapiro1965} \todo{did/did not} indicate a
departure from normality (\todo{$W$, $p$}); accordingly a
\todo{paired $t$-test / Wilcoxon signed-rank / one-way ANOVA /
Mann--Whitney} was used, giving \todo{statistic} $= \todo{}$,
$p = \todo{}$, effect size $= \todo{}$
\citep{mann1947, cohen1988}. \todo{If a between-subjects factor other than
pre/post is used (e.g., high vs.\ low engagement), define it explicitly
here and justify the grouping.}

\section{Secondary analyses}
\label{sec:results-secondary}

\todo{Report whichever of the analyses in \cref{tab:analyses} were run,
each with its figure/table. At minimum, analyses 1 (intervention
completion) and 6 (self-report vs.\ detected affect) are expected given
the primary questions; analyses 2, 3, and 7 are high-value if time
permits.}

\medskip
\noindent The next chapter interprets these results against the two
research questions and the threats to their validity.
